\subsection{Burnside's Lemma}

Burnside's Lemma (also called the Cauchy-Frobenius Lemma) is used to count the number of distinct objects under group symmetries.

\subsubsection{Statement}

Let $G$ be a group of symmetries acting on a set $X$. The number of distinct objects (orbits) is:

\begin{equation}
    |X/G| = \frac{1}{|G|} \sum_{g \in G} |X^g|
\end{equation}

where $|X^g|$ is the number of elements in $X$ that are fixed by the symmetry $g$.

\subsubsection{Interpretation}

The formula averages the number of configurations fixed by each symmetry operation over all symmetries in the group.

\subsubsection{Classic Application: Necklace Problem}

Count the number of distinct necklaces with $n$ beads and $m$ colors.

\textbf{Symmetry group:} Rotations and reflections (dihedral group $D_n$ with $2n$ elements)

For rotations by $k$ positions ($k = 0, 1, \ldots, n-1$):
\begin{equation}
    |X^{r_k}| = m^{\gcd(n,k)}
\end{equation}

For reflections (if applicable):
\begin{itemize}
    \item If $n$ is odd: $n$ reflections through a vertex and the midpoint of the opposite edge, each fixes $m^{(n+1)/2}$ configurations
    \item If $n$ is even: $n/2$ reflections through opposite vertices fix $m^{(n+2)/2}$ configurations, and $n/2$ through edge midpoints fix $m^{n/2+1}$ configurations
\end{itemize}

\subsubsection{Rotations Only}

If considering only rotations (not reflections), the number of distinct necklaces is:

\begin{equation}
    \frac{1}{n} \sum_{k=0}^{n-1} m^{\gcd(n,k)}
\end{equation}

Using Euler's totient function, this can be simplified to:

\begin{equation}
    \frac{1}{n} \sum_{d|n} \varphi(d) \cdot m^{n/d}
\end{equation}

\subsubsection{Applications}

Burnside's Lemma is used for:

\begin{itemize}
    \item Counting distinct colorings of geometric objects
    \item Necklace and bracelet problems
    \item Graph isomorphism counting
    \item Chemical enumeration (counting distinct molecules)
    \item Pattern counting with symmetries
\end{itemize}
